% Source: 07_Chapter_1_Away_With_the_King_The_Jesus_Movement.txt
\bookchapter[Away with the King!]{CHAPTER 1}{AWAY WITH THE KING!}{The Jesus Movement}

\leadin{C}{HRISTIANITY IS THE ONLY} major religion to have as its central event the humiliation of
its God.

\begin{blockquote}
“Dear dying Lamb,” believers sing, “thy precious Blood Shall never lose its power, Till all the
ransomed Church of God Be saved to sin no more.”
\end{blockquote}

Crucifixion was a barbarous death, reserved for agitators, pirates, and slaves. Jewish law cursed
“everyone who hangs on a tree” and the Roman statesman, Cicero, warned, “Let the very name of the
cross be far, not only from the body of a Roman citizen, but even from his thoughts, his eyes, his
ears.”

Part of the victim’s punishment was to be whipped and then to carry the heavy crossbeam to the place
of his own death. When the cross was raised, a notice was pinned to it giving the culprit’s name and
crime. In Jesus’ case, INRI: Iesus Nazarenus Rex Iudaeorum (Jesus of Nazareth, King of the Jews).

Pontius Pilate, Jesus’ Roman judge, apparently intended it as a final thrust of malice aimed at the
Jews, but, like the cross itself, Jesus’ followers found a special meaning in the message.

\section{JESUS AND THE CHURCH}

Jesus was a Jew. He came from a Jewish family; he studied the Jewish scriptures; he observed the
Jewish religion. Any serious study of his life makes this so clear that many people have asked if
Jesus ever intended to create that company of followers we call the church. Albert Schweitzer, the
famous missionary to Africa, believed that Jesus was obsessed with a dream of the impending end of
the world and died to make the dream come true. Rudolf Bultmann, an influential German theologian,
taught that Jesus was a prophet who challenged people to make a radical decision for or against God.
Other Christians have held that Jesus’ kingdom was a brotherhood of love and forgiveness. If he
founded a society at all, they say, it was an invisible one, a moral or spiritual company—not an
institution with rites and creeds.

This anti-institutional view of Christianity is so widespread that we had better face the question
straightaway. Did Jesus have anything to do with the formation of the Christian church? And if he
did, how did he shape its special character?

The gospel writers picture Jesus as retracing the steps of Israel. Reminiscent of Israel, Jesus
spent time in Egypt, entered the Jordan (baptism), was tempted in the wilderness, called twelve
apostles (like twelve tribes), spoke God’s word like Moses (Sermon on the Mount), preached five
sermons (compare the Pentateuch) in Matthew, performed mighty deeds of deliverance (signs, wonders,
and exorcisms), and confronted imperial powers. Where Israel had failed, Jesus had been a faithful
Son. His followers were to take up the task of being God’s servant people. He worked with a faithful
band of disciples, he taught them about life in what he called “the kingdom of God,” and he
introduced them to the new covenant that bound them together in forgiveness and love.

Granted, that simple company lacked many of the laws, officials, ceremonies, and beliefs of later
Christendom, but it was a society apart. Jesus made a persistent point about the special kind of
life that separated the kingdom of God from rival authorities among men. Little by little his
disciples came to see that following him meant saying no to the other voices calling for their
loyalties. In one sense that was the birth of the Jesus movement. And in that sense, at least, Jesus
“founded” the church.

\section{PALESTINE IN JESUS’ DAY}

During the days of Jesus, Palestine never lacked for loyalties. It was a crossroads of culture and
peoples. Its two million or more people, ruled by Rome, were divided by region, religion, and
politics. “In a day’s journey a man could travel from rural villages where farmers tilled their
fields with primitive plows to bustling cities where men enjoyed the comforts of Roman civilization.
In the Holy City of Jerusalem, Jewish priests offered sacrifices to the Lord of Israel, while at
Sebaste, only thirty miles away, pagan priests held rites in honor of the Roman god Jupiter.”

The Jews, who represented only half the population, despised their foreign overlords and deeply
resented the signs of pagan culture in their ancient homeland. The Romans were not just another in a
long series of alien conquerors. They were representatives of a hated way of life. Their imperial
reign brought to Palestine the Hellenistic (Greek) culture that the Syrians had tried to impose
forcibly on the Jews over a century before. All the children of Abraham despised their overlords;
they simply disagreed about how to resist them.

Centuries earlier the prophets of Israel had promised a day when the Lord would deliver his people
from their pagan rulers and establish his kingdom over the whole earth. On that day, they said, he
would send an anointed ruler, a messiah, to bring an end to the corrupt world of the present and
replace it with an eternal paradise. He would raise the dead and judge their actions in this world.
The wicked would be punished, but the righteous would be rewarded with eternal life in the kingdom
of God.

According to the book of Daniel and other popular Jewish writings, the Lord’s kingdom would be
established only after a final cosmic struggle between the forces of evil, led by Satan, and the
forces of good, led by the Lord. It would end with the destruction of the existing world order and
the creation of a kingdom without end (Dan. 7:13–22). This belief, along with ideas about the
resurrection of the dead and the last judgment, was in Jesus’ day very much a part of popular Jewish
faith.

Out of the distaste for life under the Romans, several factions arose among the Jews, each
interpreting the crisis in a different way. The Jesus movement was one of them.

One group, the Pharisees, emphasized those Jewish traditions and practices that set them apart from
pagan culture. Their name means separated ones, and they prided themselves on their strict
observance of every detail of the Jewish law and their extreme intolerance of people whom they
considered ritually unclean. This piety and patriotism won respect among the people.

On the other hand, some Jews found Roman rule a distinct advantage. Among them were members of
Jerusalem’s aristocracy. From this small group of wealthy, pedigreed families came the high priest
and the lesser priests who controlled the temple. Many of them enjoyed the sophisticated manners and
fashions of Greco-Roman culture. Some even took Greek names. Their interests were represented by the
conservative political group known as the Sadducees. At the time of Jesus, these men still
controlled the high Jewish council, or Sanhedrin, but they had less influence among the common
people. Another party, the Zealots, were bent on armed resistance to all Romans in the fatherland.
They looked back two centuries to the glorious days of the Maccabees when religious zeal combined
with a ready sword to overthrow the pagan Greek overlords. Thus the hills of Galilee often concealed
a number of guerrilla bands ready to ignite a revolt or destroy some symbol of Roman authority in
Palestine.

Finally came the Essenes, who had little or no interest in politics or in warfare. Instead, they
withdrew in protest to the Judean wilderness, believing the temple of Judaism to be hopelessly
compromised. There, in isolated monastic communities, they studied the Scriptures and prepared
themselves for the Lord’s kingdom, which they believed would dawn at any moment. Scholars typically
identify the Essenes as the occupants of the Quran community who copied ancient manuscripts and
wrote commentaries. These documents, called the Dead Sea Scrolls, were discovered in 1946.

Jesus had to call for the loyalty of his followers without confusing the purpose of his mission with
the objectives of these other parties among the Jews. It was a tough assignment.

\begin{sidebar}{Judaism Now and Then}
The Pharisees’ version of first-century Judaism has survived and evolved to become the Rabbinic
Judaism we know today. The destruction of the temple in AD 70 changed the character of Judaism. The
Roman resolve to quash revolutionary movements made the Zealots (a protest seeking political
revolution) and the Essenes (a protest seeking purity) impracticable.

The Sadducees were the aristocratic power brokers who operated the temple, and they saw their vision
of Judaism vanish with the temple’s destruction. The destiny of Judaism was left to the “book
people” (the Pharisees), who sought to direct the entirety of their lives by the instructions of the
Hebrew Scriptures. The Pharisees’ book-oriented version survived once the temple was destroyed and
Roman power silenced revolutionary voices.
\end{sidebar}

\section{JESUS’ MINISTRY}

Jesus chose to begin by recognizing a new movement in the Judean wilderness led by a prophet named
John. The ford of the Jordan, just north of the Dead Sea, was one of the busiest parts of the whole
region, so John the Baptist got the crowds he wanted to hear him. Wearing a garment of camel’s hair,
his eyes ablaze, he stood on the riverbank and warned all who passed by to repent of their sins and
prepare for the coming day of judgment by receiving baptism in the Jordan. Israel first entered the
land by crossing the Jordan; Jesus began his ministry at this pivotal place.

Many thought John was the promised Messiah, but he vehemently denied any such role. He explained his
mission in the words of the prophet Isaiah: “The voice of one crying in the wilderness: Prepare the
way of the Lord, make his paths straight” (Matt. 3:3). He was, he claimed, only the forerunner of
the Messiah. “I baptize you with water” he said, “but~.~.~. he will baptize you with the Holy Spirit
and with fire” (Luke 3:16).

John’s call to repentance and righteousness drew Jesus to the Jordan. He found in John’s message the
truth of God, so “to fulfill all righteousness” he submitted to John’s baptism and soon afterward
began his own mission, proclaiming, “The time is fulfilled, and the kingdom of God is at hand;
repent, and believe in the gospel” (Mark 1:15).

Jesus, however, rather than remaining in the desert, chose to begin his missionv in Galilee, a land
of gentle hills and warm, green valleys. During those early weeks and months he traveled from
village to village throughout Galilee, preaching in synagogues in the evening and on the Sabbath.
Carrying a bundle of bread, a wineskin, and a walking stick, he hiked along the dusty highways. He
probably dressed as any other traveler, in a rough linen tunic covered by a heavier red or blue
mantle.

On a typical day Jesus would set out at dawn and walk mile after mile. Toward sunset he would enter
a village and proceed to its synagogue. As one popular history puts it, “There he probably received
a warm welcome from the townspeople, who often had no resident rabbi and relied on the services of
wandering teachers like Jesus. When the lamps had been lit and the men of the village had taken
their places, Jesus would seat himself on the raised central platform” and begin reading a passage
from the sacred Scriptures. In a clear, forceful voice he would announce the fulfillment of some
prophecy or relate some parable.

The main theme of Jesus’ teaching was the kingdom of God. What did he mean by that? Did he believe
in a dramatic intervention of God in the history of the world? Or did he mean that the kingdom is
already here in some sense? He probably meant both. The two can be reconciled if we recognize that
the phrase stands for the sovereignty of a personal and gracious God, not a geographical or local
realm.

Jesus taught that the rule of God was already present in saving power in his own person. And he
offered proof of the point. His miracles of healing were apparently not just marvels; they were
signs, the powers of the age to come already manifest in the present age. “But if it is by the
finger of God that I cast out demons,” he once said, “then the kingdom of God has come upon you”
(Luke 11:20). Yet he feared that his cures would be misinterpreted, that people would see him as
just another magician, and he often cautioned those he healed to be silent.

Of course, the news spread, and before long people in every town and village in Galilee were talking
excitedly of the new wonder-worker who could cure the blind, the lame, and the sick with the power
of his voice and the mere touch of his strong carpenter’s hands. Soon large crowds gathered wherever
he spoke.

Jesus’ growing popularity aroused controversy, especially among the Pharisees, who hated to see
people following a man who had never studied under their learned scribes. They didn’t hesitate to
question his credentials openly.

\section{JESUS’ MESSAGE}

Jesus welcomed their challenge for it gave him a chance to contrast his message of repentance and
grace with the self-righteousness of the Pharisees. On one occasion, probably as pilgrims were on
their way to Jerusalem for one of the great feasts, Jesus told about two men who went to the temple
to pray. What a striking contrast they made! One was a Pharisee; the other, surprisingly, was a
despised tax collector.

With a touch of showmanship, the holy man took his stand and prayed, “God, I thank thee that I am
not like other men, extortioners, unjust, adulterers, or even like this tax collector. I fast twice
a week, I give tithes of all that I get” (Luke 18:11–12). That, at any rate, is what he prayed to
himself, and it was not a hollow boast. Pharisees excelled in those works of righteousness—fasting
and tithing—that set them apart from wicked men. The fault of the prayer was in its spirit of
self-righteousness and its cruel contempt for others. The Pharisee alone was righteous, and all his
fellow mortals were included under one sweeping condemnation.

The tax collector believed he was religiously compromised. By working to collect taxes for the
Romans, he broke faith with his own people. Sensing his own feeble religious standing, he stood at a
distance, the very image of contrition. His eyes were downcast, his head bowed in guilt. His prayer
was a sob of remorse, a cry for mercy: “God, be merciful to me a sinner!”

“I tell you,” said Jesus, “this man went down to his house justified rather than the other” (Luke
18:14). The contrast between the piety of the Pharisees and the attitude of the Jesus movement could
hardly be greater. One was based on the observance of the hundreds of religious laws of the Jews;
the other rested upon a denial of self-righteousness and a trust in the mercy of God.

Out of his hundreds of followers Jesus called a handful to travel with him full time. They came to
be called apostles, meaning sent ones. At first they were a rather motley group, twelve in all,
drawn from fishing boats and tax tables, but their loyalty to Jesus was strong.

So for them Jesus drew the distinction between his kingdom and the kingdoms of the world. His
followers, he said, represented another type of society and another type of greatness. In the
kingdoms of this world, powerful leaders lord it over others; but God’s kingdom is governed in a
wholly different way, by love and service.

“Fear not,” he told them, “it is your Father’s good pleasure to give you the kingdom” (12:32).

The highwater mark of Jesus’ popularity came about a year before his arrest in Jerusalem. After he
fed over five thousand Passover pilgrims on a grassy hillside in Galilee, many of his disciples
tried to proclaim him king. Jesus knew, however, that they had no idea of God’s unfolding plan for
his life—and death. So he fled to the hills with a committed few.

Jesus knew that he had a unique role in God’s plan of redemption, but he feared the traditional
titles for a messianic redeemer. Crowds were too likely to misunderstand them. The picture that
appears in his teaching of the twelve is along the lines of Isaiah’s portrait of the Suffering
Servant, “despised and rejected by men~.~.~. with his stripes we are healed” (Isa. 53:3, 5), and the
image of Zechariah’s predicted king who would be “humble and riding on an ass” (Zech. 9:9).

\section{THE LAST WEEK}

Apparently with these prophetic portraits in mind on the Sunday before his last Passover, Jesus rode
into Jerusalem on a donkey in fulfillment of Zechariah’s prediction. Crowds threw palm branches in
his path and shouted, “Hosanna! Blessed is he who comes in the name of the Lord.”

This seems to be the only occasion when Jesus openly identified himself with the Messiah of Jewish
prophecies. He apparently intended to challenge the Jerusalem authorities to make up their minds:
Would they or would they not accept the rule of his kingdom? The Holy City was stirred, asking, “Who
is this?”

The next day Jesus led a procession through Jerusalem’s teeming, narrow streets, to the temple.
There, in an act of protest reminiscent of the Old Testament prophets, he “entered the temple of God
and drove out all who sold and bought in the temple, and he overturned the tables of the
money-changers, and the seats of those who sold pigeons. He said to them, ‘It is written, “My house
shall be called a house of prayer”; but you make it a den of robbers’\,” (Matt. 21:12–13).

News of this dramatic event quickly swept through Jerusalem and people began flocking to the temple,
hoping to catch a glimpse of Jesus. Rumors spread of the appearance of the Messiah and the imminent
destruction of the temple.

Such talk of a messiah alarmed the temple authorities. What if this Galilean were to ignite another
revolt against the Roman government? Yet they were hesitant to arrest him for fear of provoking a
riot.

A man like Jesus presented a real danger to the Sadducees, because they held their privileged
position with the support of the Roman authorities. Anyone who aroused talk of a messiah undermined
the people’s allegiance to the established political order and endangered the relationship the
Sadducees had with the Romans. Such a man, they concluded, had to be silenced before he sparked an
uprising, which the Romans would crush with characteristic brutality. If that happened, the
Sadducees stood to lose their privileges.

Thus their common fear of Jesus brought about an unusual alliance between the Sadducees and their
rivals, the Pharisees. Jesus, who openly violated the Sabbath laws and questioned the validity of
other laws, seemed to be undermining the authority of the Jewish religion. For their separate
reasons, both parties saw this self-styled prophet from Galilee as a dangerous enemy, and together
they concluded that he should be brought to trial and condemned to death.

The temple authorities found their opportunity among Jesus’ closest followers. With the aid of Judas
from Iscariot, one of the Twelve, they could arrest Jesus secretly without provoking a riot; so
“they paid him thirty pieces of silver,” nearly four months’ wages for a skilled worker, providing
he would lead them to Jesus.

\section{THE NEW COVENANT}

“The next day was the first day of the Jewish Passover, and Jesus and his disciples prepared for the
ritual dinner that evening. At sundown they gathered secretly at the appointed place. Their mood was
solemn as they ate the meal, commemorating the Exodus of the Jews from Egypt. Reclining on couches
arranged around a low table, they drank wine and ate the bitter herbs and unleavened bread.” Toward
the end of the meal, Jesus took a piece of bread, gave thanks to God, broke it, and said, “This is
my body which is given for you. Do this in remembrance of me” (Luke 22:19). In the same way he took
a cup, saying, “This cup is the new covenant in my blood. Do this, as often as you drink it, in
remembrance of me” (1 Cor. 11:25).

What did Jesus mean by this “new covenant”? The background, to be sure, was the Exodus from Egypt
and the formation of Israel as a nation at Mount Sinai. But Jesus had in mind more than this
reminder of the obvious.

He spoke of the new covenant in his own blood. His words were an echo of the prophet Jeremiah who
had promised a day when the covenant on tablets of stone would be replaced by a covenant written on
the hearts of men: “This is the covenant which I will make with the house of Israel after those
days, says the LORD: I will put my law within them, and I will write it upon their hearts; and I
will be their God, and they shall be my people~.~.~. for I will forgive their iniquity, and I will
remember their sin no more” (Jer. 31:33–34).

The time of the new covenant, said Jesus, has come. A new people of God, enjoying the forgiveness of
sins, is now possible through the shedding of his own blood.

At that moment the disciples were undoubtedly as puzzled by his words as by his actions. But in a
matter of weeks they would see all these final hours in a new, revealing light.

After the meal Jesus led the disciples to a familiar meeting place at the foot of the Mount of
Olives, an olive grove known as Gethsemane. There was a full moon, and the grove was bathed in soft
light. While the disciples slept, Jesus withdrew to pray: “My Father, if it be possible, let this
cup pass from me; nevertheless, not as I will, but as thou wilt” (Matt. 26:39).

After renewing his commitment to God in prayer, Jesus aroused his sleeping disciples. “While he was
still speaking, Judas came~.~.~. and with him a great crowd with swords and clubs, from the chief
priests and the elders” (Matt. 26:47). They seized Jesus and dragged him away to the palace of
Caiaphas, the high priest, in the western section of Jerusalem.

\section{THE TRIAL AND DEATH}

Inside the splendidly appointed mansion, the Sanhedrin swept aside all tokens of justice and hastily
secured two witnesses who testified against Jesus. The court charged him with blasphemy and voted to
put him to death, but for that they were forced to turn to a despised Roman.

When the first rays of light appeared, the Jewish authorities led Jesus out of Caiaphas’s palace and
through the streets to the Antonia, a palace-fortress where the Roman governor, Pontius Pilate, was
staying during the Passover. Since the Sanhedrin was not empowered to carry out the death sentence,
the members had to present their case against Jesus to Pilate.

“A messenger entered the sumptuous chambers of the Antonia to summon Pilate, while the council
members and their prisoner waited below in the paved courtyard of the fortress. A few minutes later
the governor appeared. He wore a red toga draped over a white tunic in customary Roman fashion, the
distinctive mark of a Roman citizen.”

After asking about their purpose, the Roman governor pondered the situation. It seemed to him that
the chief priests had approached him to settle a petty religious dispute, and to convict Jesus
during the festival could surely spark at least a minor uprising. Yet if he ignored their
accusations and this Galilean eventually proved to be a traitor to Rome, his own position would be
endangered. Meanwhile a belligerent crowd had gathered outside the Antonia, clamoring for Pilate’s
decision. Fearful of offending Caesar, Pilate delivered Jesus to his soldiers for crucifixion.

When the execution party reached a hill outside Jerusalem called Golgotha, the soldiers stripped the
clothes from Jesus and divided the garments among themselves as the crosses were assembled. “Each
prisoner was then placed on his cross. Jesus suffered in silence as the soldiers nailed his wrists
to the crosspiece with large iron spikes and drove another spike through both ankles. As they lifted
his cross upright, his weight was supported by a peg jutting out from the cross between his legs.”
Then the soldiers fastened over the cross that sign describing his crime: “Jesus of Nazareth, King
of the Jews.”

“It was a slow and painful death. Jesus hung there helplessly for long hours as the hot sun beat
down on his body and insects buzzed about his limbs. Curious passersby paused to watch his agony and
to read the sign. Gradually he weakened, his body tortured by muscle cramps, hunger, and thirst.” A
small group of his despairing followers watched in silence as his life slipped away, a strange and
revealing prelude to the history of Christianity.

As he grew weaker and weaker Jesus cried out, “It is finished,” and yielded up his spirit. Within
hours a friend, Joseph from Arimathea, carried Jesus’ body into his own garden. There he had a tomb
hewn out of a large rock. Inside near the rear of the tomb was a couch, also of stone, and Joseph
gently placed the corpse upon it. Then he rolled a heavy stone across the entrance and went home.

\readinghead{Suggestions for Further Reading}
\begin{readinglist}
\entry Blomberg, Craig L. Jesus and the Gospels: An Introduction and Survey. Nashville: Broadman and
  Holman, 1997.
\entry Drane, John. Jesus and the Four Gospels. New York: Harper and Row, 1979.
\entry Reader’s Digest Association. Great People of the Bible and How They Lived. Pleasantville: The
  Reader’s Digest Association, 1974.
\entry *Strauss, Mark. Four Gospels, One Jesus: A Survey of Jesus and the Gospels. Grand Rapids:
  Zondervan, 2007.
\entry *Wright, Christopher. The Mission of God’s People: A Biblical Theology of the Church’s
  Mission. Grand Rapids: Zondervan, 2010.
\entry Wright, N. T. The New Testament and the People of God. Minneapolis: Fortress, 1992.
\entry Wright, N. T. Jesus and the Victory of God: The Resurrection of the Son of God. Minneapolis:
  Fortress, 1996.
\end{readinglist}
